\section{Reproducibilidad}
\label{ap:repro}
\begin{itemize}
\item \textbf{Datos:} \url{https://github.com/Juferoga/linux-kernel-commits-jit}, \emph{release} v1.0 (\texttt{kernel-commits-2019-2026.parquet}, 31 MB, SHA-256 \texttt{ac6858f4\ldots45fe9e}) \cite{DatosKernel2026}.
\item \texttt{...-proyecto-e3-cuaderno.ipynb}: cuaderno autocontenido y ejecutado, con salidas visibles. Lee la copia incluida o descarga y verifica los datos, reserva la prueba y reproduce las Tablas~\ref{tab:particion}, \ref{tab:pipeline}, \ref{tab:espacio} y \ref{tab:referencias}. Al ejecutar las celdas, escribe el registro de experimentos (\texttt{registro\_experimentos\_e3.csv}).
\item \texttt{...-proyecto-e2-cuaderno.ipynb}: exploración del Entregable 2 (Tablas~\ref{tab:d1d7} y \ref{tab:desc}, Figs.~\ref{fig:objetivo} y \ref{fig:corr}).
\item Entorno: Python \cifEtresPython{}, pandas \cifEtresPandas{}, NumPy \cifEtresNumpy{} y scikit-learn \cifEtresSklearn{} (\texttt{requirements.txt}). Semillas: 11, 23, 37, 41 y 53.
\end{itemize}

\textbf{Tabla A1. Diccionario de variables.} Está en la hoja \emph{diccionario} del libro de Excel del anexo y en el README del repositorio de datos.
